\documentclass{article}

\usepackage{xeCJK}
% \setCJKmainfont{WenQuanYi Micro Hei}

\usepackage{geometry}
\geometry{a4paper,scale=0.8}
% \geometry{a4paper,left=2cm,right=2cm,top=1cm,bottom=1cm}

\usepackage{amsmath}

\usepackage{cleveref}

\usepackage{indentfirst}

\title{Bellman Equation Deduction}
\author{dreamanddead@foxmail.com}
\date{\today}

\begin{document}
 
\maketitle

%\begin{abstract}
% abstract
%\end{abstract}

% \newpage

\section{prequisite}

\subsection{条件概率和独立性}

条件概率是概率论是最为基础的概念。在事件 B 发生的条件下，A 事件发生的概率等于 A B 事件同时发生的概率除以 B 事件发生的概率。

$$
P(A|B) = P(AB) / P(B)
$$

移项得到

\begin{equation} \label{cond}
 P(AB) = P(A|B)P(B)
\end{equation}

如果 A B 两个事件相互独立，彼此没有依赖关系，则 A B 同时发生的概率等于 A 事件发生的概率乘以 B 事件发生的概率。

$$
P(AB) = P(A)P(B)
$$

在相互独立的前提下，综合可知

\begin{equation} \label{ind}
 P(A|B) = P(A)
\end{equation}


\subsection{全概率公式}

若事件 $B_1, B_2, ..., B_n$ 是全空间 $\Omega$ 的一个划分，
 
$$
\sum_{i=1}^n P(B_i) = 1
$$
 
则

$$
P(A) = \sum_{i=1}^n P(AB_i)
$$
 
进一步地，根据 \cref{cond} 可得
 
\begin{equation} \label{full}
 P(A) = \sum_{i=1}^n P(A | B_i) P(B_i)
\end{equation}

\cref{full} 是理解之后公式变化的要点。
 

\section{value function}

RL 的基础思想很简洁，Agent 不断和 Env 相交互，以最大化收益为目标。

在 MDP 框架下，Agent 选择 Action 影响 Env，Env 反馈 State 和 Reward。

$$
S_0, A_0, R_1, S_1, A_1, R_2, S_2, A_2, ...
$$

Env 的行为完全由 $p(s', r|s, a)$ 来决定，

$$
\sum_{s'} \sum_r p(s', r|s, a) = 1
$$


Agent 在 t 时刻之后所能得到的全部收益为

\begin{equation}
 G_t = \sum_{k=0}^{\infty} \gamma^k R_{k+t+1}
\end{equation}

当然 $G_t$ 不是孤立的存在，因为后续的收益和 agent 当前所处的状态以及后续使用的策略（policy）相关。

定义 state value function，衡量 agent 在 t 时刻处于 s 状态，后续使用策略 $\pi$ 可得到的累计收益。

 $$
 V_\pi(s) = E_\pi[G_t | S_t = s]
 $$
 
进一步地，假如 agent 已经 在 t 时刻已经采取某一动作 a，定义 state action value function 为在此之后可得到的累计收益。

 $$
 Q_\pi(s, a) = E_\pi[G_t| S_t = s, A_t = a]
 $$

事实上，由于行为和反馈不断循环的特性，$V_\pi(s)$ 和 $Q_\pi(s, a)$ 存在着一定的关系。
 
\begin{align}
 V_\pi(s) &= E_\pi[G_t | S_t = s] \nonumber \\
          &= \sum_a \pi(a|s) E_\pi[G_t| S_t = s, A_t = a] \nonumber \\
          &= \sum_a \pi(a|s) Q_\pi(s, a) \label{vq}
\end{align}

 
推导过程利用了 \cref{full}，其中
 
 $$
 \sum_a \pi(a|s) = 1
 $$
 
 
至于 Q 则复杂一点，
 
\begin{align}
Q_\pi(s, a) &= E_\pi[G_t| S_t = s, A_t = a] \nonumber \\
            &= E_\pi[R_{t+1} + \gamma G_{t+1}| S_t = s, A_t = a] \nonumber \\
            &= E_\pi[R_{t+1} | S_t = s, A_t = a] + \gamma E_\pi[G_{t+1}| S_t = s, A_t = a] \nonumber
\end{align}
 
分成两部分来看，前一部分，
 
 $$
 E_\pi[R_{t+1}| S_t = s, A_t = a] = \sum_r p(r | s, a) r
 $$
 
而
 
 $$
 p(r | s, a) = \sum_{s'} p(s', r | s, a)
 $$
 
（依旧是 \cref{full}，看出来了吗）
 
综合可得
 
 $$
 E_\pi[R_{t+1}| S_t = s, A_t = a] = \sum_r p(r | s, a) r = \sum_r \sum_{s'} p(s', r | s, a) r
 $$
 
再来看另一项
 
 $$
 E_\pi[G_{t+1}| S_t = s, A_t = a] = \sum_{s'} p(s' | s, a) E_\pi [G_{t+1} | S_{t+1} = s', S_t = s, A_t = a ]
 $$
 
（依旧是 \cref{full}）
 
其中
 
 $$
 p(s' | s, a) = \sum_r p(s', r | s, a)
 $$

（依旧是 \cref{full}）

根据 MDP 的定义， $S_t A_t$ 只影响 $R_{t+1}$，和 $R_{t+2}$ 并不相关（相互独立），回忆 \cref{ind}，综合可知
 
 $$
 E_\pi[G_{t+1}| S_t = s, A_t = a] = \sum_{s'} \sum_r p(s', r | s, a) E_\pi [G_{t+1} | S_{t+1} = s'] = \sum_{s'} \sum_r p(s', r | s, a) V_\pi(s')
 $$
 
 
 \begin{equation} \label{qv}
   Q_\pi(s, a) = E_\pi[G_t| S_t = s, A_t = a] = \sum_{s'} \sum_r p(s', r | s, a)[r + \gamma V_\pi(s')]
 \end{equation}

  
\section{recursive format}
 
利用 \cref{vq} 和 \cref{qv}，进一步推出各自 value function 的递归形式，

$$
 V_\pi(s) = \sum_a \pi(a|s) Q_\pi(s, a) = \sum_a \pi(a|s) \sum_{s'} \sum_r p(s', r | s, a)[r + \gamma V_\pi(s')]
$$
 
 
$$
   Q_\pi(s, a) = \sum_{s'} \sum_r p(s', r | s, a)[r + \gamma V_\pi(s')] = \sum_{s'} \sum_r p(s', r | s, a)[r + \gamma \sum_{a'} \pi(a'|s') Q_\pi(s', a')] 
$$
 
\end{document}














